\chapter{ALL'ARREMBAGGIO}\label{allarrembaggio}

\hfill\break

In piedi sulla porta della cabina di pilotaggio, mentre Desai faceva
girare i motori per testare il decollo, esaminai i volti della nostra
task force, o qualsiasi cosa fossimo. Avevamo davvero un gruppo
arcobaleno internazionale e multirazziale. Oltre a me, c'erano nove
elementi dell'esercito degli Stati Uniti, ovvero il maggiore Simms, un
sergente, tre specialisti e quattro soldati semplici. Il sergente scelto
Adams era il nostro solo Marine degli Stati Uniti, e avevamo un sergente
dell'Aeronautica che si occupava della manutenzione delle Poiane su
Paradiso: pensai che ci sarebbe potuto tornare utile se qualcosa su una
nave avesse avuto bisogno di riparazioni. Quattro cinesi, tra cui il
tenente colonnello Chang. Tre elementi dell'esercito indiano, compreso
il nostro unico pilota. Quattro inglesi, tra cui un sergente; uno dei
soldati semplici britannici aveva preso le lezioni di volo in un aereo
monomotore, perciò, a rigor di logica, lo affidai come copilota al
capitano Desai. Sembrava terrorizzato quando prese posto sul sedile di
destra nella cabina di pilotaggio. Ah, e un soldato dell'esercito
francese, un certo tenente Renée Giraud. Era assegnato a un commando di
paracadutisti ed era una sorta di versione francese di un ranger
dell'esercito. Avevo lavorato con le forze speciali francesi in Nigeria,
e quei ragazzi sono davvero tosti, ero felice di avere Renée in squadra.
Quando si arruolò, mi disse di non essere sicuro che il mio discorso non
fosse una stronzata, ma aveva voglia di azione e su Paradiso non ce
n'era. Apprezzai la sincerità.

Ventiquattro persone, riunite in fretta e furia. Cinque donne,
diciannove uomini. Cinque ufficiali, cinque sergenti, quattordici
reclute. Cinque nazionalità: il che però, in alcuni casi, era una
definizione un po' vaga. Uno dei nostri specialisti dell'esercito degli
Stati Uniti era un indiano d'America di nome Randy Putri e, guardandolo,
si sarebbe detto che appartenesse al contingente dell'esercito indiano,
ma non spiccicava una parola di hindi e parlava con l'accento cajun con
cui era cresciuto a New Orleans. Il nostro sergente dell'Aeronautica
degli Stati Uniti si chiamava Chung ed era cinoamericana. Chung parlava
un po' di cinese, ma il poco che sapeva era cantonese, non mandarino,
perciò non poteva comunicare con i cittadini cinesi meglio di quanto non
potessi fare io. Il sergente Reginald Thompson dell'esercito britannico
aveva la pelle scura dei suoi nonni kenioti, ma quando apriva bocca
parlava come Sherlock Holmes, o come la crème de la crème fighetta di
qualche programma della Bbc. Uno degli altri inglesi aveva un accento
così stretto che all'inizio non ero certo che parlasse inglese; il loro
slang è del tutto incomprensibile. Il capitano Desai, per qualche strana
ragione, sembrava capire il suo accento pseudo-inglese, avrebbe potuto
fare da interprete se necessario. A proposito di interpreti, il tenente
colonnello Chang e uno degli altri cinesi parlavano inglese, i due
restanti si sarebbero affidati ai traduttori dei loro zPhone. Skippy,
ovvio, parlava alla perfezione ogni lingua umana, quella piccola molesta
lattina di birra.

Ventiquattro persone, diverse per nazionalità, generi, specialità ed
esperienze, e dovevo trasformarle in una forza di combattimento
efficace, in fretta. Senza sapere con esattezza quale fosse la nostra
missione.

La cosa più preoccupante era che, tra i ventiquattro volontari, non
c'era neanche un medico. E le forniture mediche che avevamo caricato a
bordo erano piuttosto basiche. Le cose sarebbero potute andare a rotoli
molto in fretta e non ci sarebbe stato molto da fare per i feriti.

\hfill\break

Il capitano Desai ci fece decollare sani e salvi, dopo di che il pilota
automatico prese il comando del Dodo, effettuando una cabrata per uscire
dall'atmosfera. La gravità calò in modo graduale e le persone si
adattarono, assumendo farmaci contro la nausea al bisogno. Skippy,
voglio dire Hacker, disse più volte a Desai di riprogrammare il pilota
automatico, per evitare di volare vicino alle astronavi ruhar, c'era un
sacco di traffico attorno a Paradiso. Stando agli schermi della cabina
di pilotaggio e a quello che Skippy mi sussurrava all'orecchio,
scivolammo lisci come l'olio attraverso uno squadrone di fregate ruhar,
senza che ci indagassero, sfidassero o addirittura notassero; Skippy si
era infiltrato nei loro sensori e aveva fatto in modo che ci
ignorassero. Comunque l'avesse fatto, funzionò, in modo impressionante.
Quando Desai riferì che avevamo superato la velocità di fuga, lasciato
l'orbita ed eravamo soli al sicuro nello spazio interplanetario, decisi
che era il momento di parlare all'equipaggio, che mi guardava con ansia.
Raggiunsi galleggiando la porta della cabina di pilotaggio, così che
tutti potessero sentirmi.

«È il momento della trasparenza totale, gente. Non potevo informarvi
sulla nostra missione finché non avessimo lasciato l'orbita, perché non
possiamo rischiare che i kristang o i ruhar sappiano cosa stiamo
facendo. La verità è che», lanciai a Adams uno sguardo colpevole, «non
sono stato del tutto onesto con voi, per ragioni di sicurezza operativa.
Hacker non è il nome in codice di una cyber-squadra dell'Unef. Hacker è,
ehm, qui, ve lo mostrerò.» Tirai fuori Skippy dal mio zaino e commisi
l'errore di tenerlo in vista senza prima guardarlo.

Grande errore.

Quel piccolo spocchioso idiota aveva trasformato la sua superficie di
cromo lucido in un'imitazione a colori di una lattina di Bud Light Lime.
Non sapevo che potesse cambiare aspetto! Non fosse che mancava una
linguetta sopra il coperchio e che il suo fondo era quasi piatto,
avrebbe potuto ingannarmi.

Le facce dei membri dell'equipaggio davanti a me passarono in un batter
d'occhio dallo stupore al divertimento all'orrore e io li guardai, a mia
volta con orrore, mentre realizzavano di essere andati nello spazio con
un pazzo delirante. Un pazzo che aveva un amico immaginario a forma di
lattina di birra. «No no no!» Agitai la mano sinistra, mentre con la
destra scuotevo Skippy con gesto frenetico. «Skippy, cazzo, non è
divertente!» Adams si stava preparando ad attraversare in volo lo
scompartimento per raggiungermi e la sua faccia era tutt'altro che
amichevole. «Skippy, cazzo, piccolo stronzo!»

«Ah ah ah!» Skippy rise da matti e mutò di nuovo la sua superficie, ora
era una lattina magica di Coors Light. «Heeey hooo²³, gente! Sono
Skippy, l'onnipotente. Ah ah ah! Oh, cavolo, avresti dovuto vedere la
tua faccia, colonnello Joe. Impagabile.»

Feci un paio di profondi respiri, guardando Adams che stava ancora
cercando di decidere se soffocarmi. «Con tutta la sua intelligenza,
Skippy è uno stronzo al 100\%.»

«Vero, vero», ammise Skippy. «Così va meglio?», tornò al monotono cromo.

Chang ringhiò, letteralmente: «Chi cazzo è Skippy?».

«Ha uno spiegone da fare», disse Adams, senza traccia di ironia.

Notai che nessuno di loro aveva aggiunto ``signore''.

Feci un profondo respiro: «Skippy è questo, il suo nome in codice è
Hacker», indicai il suo lucido essere. «Questa cosa qui, che sembra una
lattina di birra, è un'intelligenza artificiale\ldots»

«Una mia piccolissima manifestazione nello spazio-tempo\ldots»,
m'interruppe Skippy.

«Skippy, vuoi stare zitto un minuto? Un'intelligenza artificiale
costruita dagli esseri che conosciamo come gli Anziani o ``i Primi'',
comunque si chiamino. La super-civiltà che abitava prima dei rindhalu la
galassia, gli esseri che hanno costruito i wormhole. Hanno abbandonato
la loro esistenza fisica molto tempo fa e si sono lasciati alle spalle
le Sentinelle. E hanno lasciato questa intelligenza artificiale che ha
milioni di anni. Un'intelligenza artificiale d'immensa potenza: ha
sbloccato i comandi su questo Dodo, ha disabilitato le armi dei ruhar,
ora ci sta celando ai sensori dei criceti e sta trasmettendo i corretti
codici Iff ai kristang. È così che saliremo a bordo di una loro nave e
la ruberemo, e faremo lo stesso con una nave madre thuranin. A quel
punto, disattiveremo il wormhole vicino alla Terra, così né i ruhar né i
kristang potranno più avere accesso alla nostra casa.»

«Porca di quella puttana», proruppe Simms incredula.

«Sì, è stata la mia reazione quando i ruhar mi hanno rinchiuso dentro un
magazzino e una lattina di birra su uno scaffale polveroso ha iniziato a
parlarmi.»

Adams non era convinta: «Signore, sto cercando di venirne a capo. Stiamo
affidando il nostro futuro e il destino dell'umanità a una lattina di
birra parlante?».

«Se la metti così\ldots»

«Una lattina di birra super-intelligente, anzi, incredibilmente,
inconcepibilmente intelligente!», protestò Skippy.

«Adams, dimentica quello che Skippy sembra a noi. Hai visto cos'è
successo al magazzino, i ruhar non hanno potuto usare le loro armi, ma
noi sì. Non ero io, non era l'Unef, era Skippy. Skippy ha hackerato
questo Dodo perché potessimo pilotarlo, di sicuro non sono stato io e
non è stata l'Unef. Siamo appena usciti dall'orbita, proprio in mezzo a
una flotta ruhar, e non ci hanno visto. Skippy ha hackerato i sensori
ruhar e ha fatto in modo che ci ignorassero. L'Unef non può fare niente
di tutto questo. Skippy è l'arma definitiva, il nostro asso nella
manica, ed è sicuro come la morte che l'umanità ha bisogno di un asso
adesso. Con Skippy, possiamo disabilitare, chiudere, spegnere, l'unico
wormhole che permette alle lucertole di accedere alla Terra.»

Adams ancora tentennava: «Come facciamo a sapere che non è un trucco
delle lucertole, che questa intelligenza artificiale non sta lavorando
per loro?».

«Adams, se ti viene in mente qualcosa, qualsiasi cosa, che le lucertole
ci avrebbero guadagnato aiutandoci a fuggire dai ruhar e a rubare uno
dei loro Dodo, per favore dimmelo. Perché a me non viene in mente
niente. Skippy costruirà la nostra fiducia nei suoi confronti nello
stesso modo in cui noi costruiremo la sua nei nostri: un'azione alla
volta. Ha fatto uscire di prigione me, te, Chang e Desai, e ora siamo
lontani dalla superficie di Paradiso. Se fossimo da soli, staremmo
ancora cercando di capire come aprire la porta di questa cosa», indicai
la cabina di pilotaggio del Dodo. «Ho anche pensato che potesse essere
un trucco dei ruhar per farci salire su una nave kristang, e che i
criceti avessero piazzato una bomba su questa cosa, questo Dodo in
particolare. Ancora una volta, non mi viene in mente un solo motivo per
cui i ruhar si sarebbero presi tutto quel disturbo. Questa guerra è
andata avanti per un tempo molto lungo senza umani coinvolti. Non hanno
bisogno di noi; le lucertole, i criceti, i thuranin, nessuno di loro ha
bisogno di noi umani primitivi per qualsiasi cosa. Forse è tutto uno
stratagemma elaborato per qualche motivo che non riusciamo a vedere. Ok,
forse lo è. O forse Skippy sta dicendo la verità e abbiamo la
possibilità di chiudere il wormhole che dà alle lucertole accesso alla
Terra. Se c'è una qualche opportunità reale, qualunque essa sia,
dobbiamo coglierla. Sapremo se possiamo fidarci di Skippy quando avremo
il controllo di una nave da guerra kristang. È semplice.» Skippy era
rimasto stranamente in silenzio per tutto il tempo in cui avevo parlato.
Parte dell'essere super-intelligente è sapere quando tenere la bocca
chiusa e lasciare che l'altro parli.

Chang, che era rimasto in silenzio, ma aveva scambiato sguardi sempre
più tesi con i suoi tre soldati cinesi, si schiarì la gola:
«Colonnello», disse con enfasi, «ci siamo offerti volontari per questa
missione, senza conoscere i dettagli. Mi aspettavo che il piano fosse
quello di attaccare i kristang. Ora ci stai dicendo che la nostra
missione è salvare il mondo? Salvare la Terra dai kristang?».

«Ehm», suonò così drammatico quando lo disse in quel modo, «sì. Sì,
stiamo per chiudere il wormhole che consente ai kristang l'accesso alla
Terra. La situazione tornerà com'era prima dello spostamento del
wormhole.» Davo per scontato che tutti avessero sentito le voci a
riguardo ormai. «La Terra sarà tutta sola in mezzo al nulla e questa
guerra andrà avanti senza umani.»

«Salvare il mondo?», ripeté Simms incredula.

«Questa è l'idea», annuii. «So che è una specie di cliché, ma\ldots»

«No, salvare il mondo mi suona bene.» Disse Simms pensosa. «Merda. È
vero?»

Chang spense il suo zPhone e disse qualcosa in mandarino ai suoi tre
soldati, poi tornò da me: «Colonnello, trovo tutto questo difficile da
credere. Tuttavia, non molto tempo fa, ero di stanza in un avamposto
dell'esercito sul confine mongolo. Ora, sono su una nave spaziale
aliena, a mille anni luce da casa. Il possibile è stato ridefinito così
tante volte che sono pronto ad accettare quasi tutto. Perciò, se questa
missione ha anche solo una piccola chance di liberare il nostro pianeta
dai kristang, faremo del nostro meglio».

Cazzo, il suo inglese era meglio del mio: «Grazie, colonnello Chang».
Esaminai le altre facce.

Giraud fece spallucce: «Uccideremo dei kristang, no? Io ci sto, come
dite voi».

Simms e Adams si scambiarono un'occhiata: «Ah, che cazzo, certo, anche
noi», dichiarò la prima. «Colonnello, ci siamo, giusto? È tutta la
verità? Niente più sorprese? Siamo alleati con una lattina di birra
parlante, per chiudere il wormhole? Ora che siamo qui, mi aspetto che la
necessità del processo di sicurezza operativa sia finita fuori dalla
finestra», guardò in modo significativo la porta della camera
d'equilibrio, «se quest'espressione si può usare nello spazio.»

Annuii: «Sapete tutto quello che so io. Non posso promettere che non ci
saranno altre sorprese, ma lo saranno anche per me».

«Perché noi, signore?», chiese il sergente Thompson nel suo inglesissimo
accento britannico. «Perché non portare questo, Skippy, al quartier
generale dell'Unef e lasciare che fossero loro a sbrigarsela? Inviando
una vera forza operative speciale, Sas e tutto il resto?»

Chang e Simms fecero all'unisono una risata nasale denigratoria.
«Sergente», disse Chang, «all'Unef si prenderebbero prima una settimana,
come minimo, per capire cosa fare. Poi, essendo avversi al rischio, è
probabile che consegnerebbero il nostro amico, l'intelligenza
artificiale, ai ruhar o ai kristang o a chiunque fosse a capo del
pianeta in quel momento, nella speranza di ottenere favori. In ogni
caso, non farebbe differenza: una volta che tutto questo si sapesse al
quartier generale dell'Unef, non ci sarebbe modo di tenerlo segreto a
lungo. I criceti e le lucertole ne verrebbero presto a conoscenza.»

«Inoltre», aggiunsi, «dobbiamo andare ora, per approfittare della
situazione. I ruhar sono impegnati a consolidare la loro presa sul
pianeta, le navi stanno saltando dentro e fuori là sopra, e c'è ancora
un residuo di forza kristang in giro. Se aspettiamo, le lucertole si
ritireranno e avremo bisogno di attaccare i ruhar per rubare una nave. I
ruhar hanno troppe navi che si sostengono a vicenda per poterne prendere
una e sgattaiolare via.»

«Giusto.» Thompson sembrava soddisfatto. «Un'altra domanda, se posso:
perché l'intelligenza artificiale si chiama Skippy?» Dal numero delle
teste che annuivano in tutto il Dodo, sembrava essere una questione
universale.

«Perché», tenni Skippy di fronte a me e lo guardai storto, «come ho
detto, è super-intelligente, super-potente e super-stronzo. Voleva
essere definito il Signore Dio Onnipotente, così l'ho chiamato Skippy,
per ricordarmi che testa di cazzo è.»

«Colpevole di tutte le accuse», disse Skippy in tono allegro.

«Non gli dispiace essere chiamato, ehm, Skippy?», chiese Adams con
franchezza.

«A me non dispiace essere chiamato Skippy, non c'è bisogno di chiedere
al colonnello Joe, sono una persona», disse Skippy. «No, non m'importa
come mi chiamate voi, branco di scimmie morsicate dalle pulci.»

«Scimmie?», chiese Putri. Lo specialista dell'esercito americano Randy
Putri, non il soldato semplice Asok Putri dell'esercito indiano. Lo so,
mi confondo anch'io. In futuro romperò la tradizione militare e li
chiamerò con il nome proprio.

«Pensa di essere gentile considerandoci scimmie», gli spiegai, «per lui
siamo tutti batteri.»

«Ah ah!», sbeffeggiò Skippy. «Vi piacerebbe essere intelligenti come i
batteri.»

«Giusto, allora», disse Thompson in tono monocorde. «Questo bastardo è
un vero stronzo.»

Adams fece la linguaccia a Skippy, un gesto che non mi sarei mai
aspettato da lei. «Qual è la prossima mossa, signore?»

Dentro di me, tirai un sospiro di sollievo: «Skippy, puoi caricare
schemi di navi kristang verosimili sugli zPhone di tutti e mostrarli
anche su questo schermo?».

«Fatto», si limitò a rispondere lui, e lo schermo sulla paratia dietro
di me si accese. «Questa è una tipica fregata kristang, il genere di
piccola nave che con molta probabilità verrà inviata a prenderci\ldots»

Un vero colonnello era responsabile del comando di una forza delle
dimensioni di una brigata. Ciò significava pianificare operazioni
offensive e difensive di diverse migliaia di soldati, compresi tutti gli
addestramenti, la logistica, le comunicazioni e il coordinamento con la
potenza aerea, l'artiglieria e altre unità nella zona. Fare ed eseguire
grandi progetti che coinvolgevano migliaia di persone ed enormi potenze
di fuoco. Un vero colonnello aveva una formazione, un'istruzione teorica
e pratica e anni di esperienza prima di assumere il comando. Io non
avevo niente di tutto ciò.

Quello che avevo era competenze acquisite con fatica in combattimenti di
piccole unità, esperienza nella boscaglia e, cosa più importante, in
villaggi e città della Nigeria settentrionale. Il furto di una nave
kristang era un combattimento di piccola unità, e pensavo che ripulire
compartimento per compartimento una nave fosse simile a combattere casa
per casa e stanza per stanza. L'esercito statunitense chiamava questo
tipo di conflitto ``operazioni militari in terreno urbano'' e le basi
militari negli Stati Uniti contenevano villaggi simulati per addestrare
le truppe a combattere in spazi tanto ristretti. Durante la guerra
fredda, questi villaggi simulati erano stati creati in modo da
assomigliare a quelli dell'Europa orientale, con cartelli sulle strade e
sugli edifici in una lingua d'impronta slava o germanica. Di recente,
l'attenzione si era concentrata sul Medio Oriente, e noi soldati avevamo
dato alle città di addestramento nomi politicamente scorretti, come
Hadjistan. Qualunque fosse il nome, o il grande scenario per cui le
città erano state costruite, avevano addestrato i soldati a sgomberare
un'area edificio per edificio e stanza per stanza; il tipo di
combattimento in cui spesso non potevi vedere a chi stavi sparando, e
chiamare dentro le navi da guerra Apache significava ritirarsi in
fretta, in modo che i missili Hellfire non colpissero le tue truppe.

Incaricai il tenente Giraud di pianificare il nostro attacco e fornire
un rapido addestramento al nostro equipaggio. Chang, Simms e Desai erano
superiori a Giraud per grado, ma Chang aveva fatto esperienza
nell'artiglieria, non nella fanteria. Simms era un ufficiale logistico e
Desai un pilota. Potreste immaginare che il sergente Adams, in quanto
Marine, fosse a rigor di logica la persona più adatta a pianificare
un'operazione d'imbarco su una nave nemica, e sarebbe stata una buona
idea, ma nella guerra anglo-americana del 1812. Il corpo dei Marine
degli Stati Uniti era un po' arrugginito su tutta la faccenda
dell'andare ``all'arrembaggio''. Restavamo io e Giraud, in quanto
ufficiali di fanteria, e Giraud era delle forze speciali francesi.
Quello che sapeva lui delle tattiche delle piccole unità e quello che
non ne sapevo io un po' mi spaventava, i ragazzi delle forze speciali
erano dei duri. Continuavo a ricordargli che la nostra forza d'assalto
non era l'élite di killer delle forze speciali cui era abituato e che
avremmo operato a gravità zero, senza addestramento e con persone che
per lo più non conoscevamo. Decidemmo che avremmo dovuto in gran parte
improvvisare il nostro attacco, perché c'erano troppe incognite. Quello
che volevo era guidare di persona una forza che s'impossessasse del
centro di controllo della navicella, che era a poppa del ponte sulla
maggior parte delle navi kristang, mentre il gruppo di Giraud si sarebbe
impossessato della sala macchine, prima che le lucertole lì potessero
danneggiare le unità di trasmissione o il reattore, o addirittura
autodistruggere la nave. Quel piano fu respinto all'unanimità da tutti,
incluso Skippy.

«Colonnello», dichiarò Adams con le braccia incrociate sul petto, in un
gesto che non era facile a gravità zero, «non può andarsene in giro di
corsa per la nave nemica. Lei è il comandante. Deve restare indietro e
comandare.»

«Adams ha ragione, colonnello Joe», mi ammonì Skippy. «Posso dirti dove
sono i tuoi e i kristang e cosa stanno facendo, e posso controllare
parti della nave, ma ho bisogno che qualcuno mi dica cosa vuoi che
faccia. E dove dovrebbero andare i tuoi. Il mio genio non comprende
coordinare una truppa di scimmie in combattimento. Sai ciò che possono
fare i tuoi e, cosa più importante, ciò che non possono fare.»

Dopo avere protestato che la mia esperienza era a livello di squadra e
gruppo di fuoco e che Giraud sarebbe stato un candidato migliore per
rimanere con Skippy e coordinare l'attacco, fui costretto a riconoscere
che avevano ragione. Alla fine, l'elemento tacito che risolse la
questione fu il fatto che mi trovavo a mio agio con Skippy e lui a suo
agio con me, e nessun altro voleva il compito di avere a che fare con
quella nervosa intelligenza artificiale aliena.

Giraud indusse il nostro equipaggio a correre con la mente attraverso un
attacco simulato su una tipica fregata kristang, usando la migliore
ipotesi di Skippy sull'ipotetica posizione dell'equipaggio kristang e su
quali sistemi a bordo della fregata avrebbe potuto controllare. Rimasi
sorpreso quando Skippy disse di non poter prendere il controllo totale
dei sistemi informatici della fregata, almeno non subito. I kristang
avevano costruito le loro navi in modo da impedire che un nemico potesse
hackerare i loro computer da remoto, dato che avevano paura che i loro
patroni, i cyborg thuranin, lo facessero. Proprio perché i thuranin
l'avevano già fatto in precedenza, e i kristang non sarebbero rimasti
vittime di quel trucco per la seconda volta o, come riferì Skippy, più
facile fosse la centesima. I kristang erano diventati più bravi a
rafforzare i loro sistemi informatici contro le intrusioni, ma il
livello della loro tecnologia era ancora tanto al di sotto di quella dei
thuranin che non erano in grado d'immaginare nemmeno una parte dei
sistemi di cui i loro patroni erano soliti servirsi. Per nostra fortuna,
Skippy superava di gran lunga i thuranin per livello tecnologico, così
come i thuranin superavano gli umani. O, come disse Skippy con grande
tatto, un albero pieno di scimmie.

La chiave di tutto il piano era che dovevamo di fatto collegare Skippy
alla nave, per stabilire una connessione fisica tra lui e il network di
quest'ultima. Una volta entrato, avrebbe avuto accesso all'intera nave e
avrebbe potuto scaricare una sub-routine di se stesso nei suoi nodi di
controllo. Ma prima dovevamo localizzare una presa a muro a bordo della
fregata, e spinottarci uno zPhone. Per fortuna, a bordo del Dodo c'erano
i cavi e i connettori necessari. Purtroppo, raggiungere la presa a muro
più vicina ai due hangar di atterraggio della fregata sarebbe stata una
gran rottura di coglioni. E il trucco di Skippy di disabilitare le armi,
come aveva fatto con i fucili ruhar, non avrebbe funzionato con i
kristang, perché realizzavano i loro fucili per lo più senza sofisticati
chip informatici, non fidandosi di quel tipo di tecnologia.

Il piano di Giraud prevedeva l'impiego dell'intera forza d'imbarco di
ventidue persone nello sforzo di stabilire una connessione fisica, dopo
di che il nostro equipaggio si sarebbe diviso in due gruppi, uno guidato
da Giraud e uno guidato da Chang. Al fianco di Giraud, designai il
sergente Thompson, al fianco di Chang il sergente Adams e il sergente
cinese. Vedevo che le persone erano molto tese, perché comprendevano più
che bene che inserire uno zPhone in una presa a muro era una questione
di vita o di morte, non solo per noi, ma in potenza per tutta l'umanità.
Non c'era bisogno di discorsi motivanti. Ordinai a tutti di mangiare e
bere qualcosa e rilassarsi. Giraud accettò, aggiungendo che voleva che
l'equipaggio memorizzasse la struttura di fregate e cacciatorpediniere
kristang di cui, per fortuna, esistevano solo due tipi principali.
«Immaginate la struttura nella vostra mente», consigliò, più e più
volte, «fino a quando non è diventata istinto, radicato nella vostra
memoria spaziale.»

«Penso che abbiamo il miglior piano possibile», dissi a Giraud, mentre
mangiavamo entrambi una barretta energetica Hooah!. L'attesa, con il
Dodo che si allontanava in volo non propulso da Paradiso, sperando di
essere imbarcato da una nave da guerra nemica, stava dando sui nervi a
tutti.

Giraud fece un'alzata di spalle esagerata: «I piani sono l'inizio,
niente di più».

«Nessun piano sopravvive al contatto con il nemico, giusto? Lo diceva
von Moltke.» Lo avevo imparato dall'esercito da qualche parte lungo la
strada.

Giraud arricciò il naso: «Von Moltke lo imparò da Napoleone. Lui
sottolineava l'importanza di avere capacità flessibili e di sfruttare le
opportunità sul campo di battaglia, piuttosto che cercare di attenersi a
piani dettagliati». Intervenne Skippy: «Nessun piano dell'Unef avrebbe
potuto prevedere questa opportunità».

\hfill\break

Non c'era posto per la privacy a bordo del Dodo, fatta eccezione per il
solo bagno a gravità zero, che era sempre in uso. Quando Skippy mi disse
che voleva parlare in modo semiprivato, avanzai verso l'angusta cabina
di pilotaggio, per galleggiare dietro il sedile destro. «Che cosa c'è,
Skippy?»

Lui mi rispose nell'auricolare. «Volevo solo dire qualcosa di carino,
per una volta, e non voglio che mi sentano, perché mi rovinerebbe la
reputazione da duro.»

Reputazione da duro? Ma Skippy dove diavolo prendeva le sue nozioni
sulla cultura umana? «Ah, certo.» Non sapevo cos'altro dire.

«La tua specie ha una straordinaria adattabilità. La capacità di
accettare nuove informazioni, nuovi concetti, di non fuggire urlando o
paralizzarsi di fronte a cambiamenti sconvolgenti, è troppo rara fra le
specie intelligenti. Mi aspettavo che ci sarebbero stati dei problemi
quando hai rivelato la verità su di me e sulla nostra missione, ma i
tuoi hanno accettato la nuova situazione in modo ammirevole e veloce.»

«Ehm.» Skippy non stava tenendo in conto che il nostro equipaggio era
formato da soldati. I soldati sono abituati al fatto che cambino loro di
continuo le carte in tavola nel bel mezzo di una missione. Una volta, in
Nigeria, ero in un Blackhawk che sfiorava le cime degli alberi alle
03:30, andavamo a fare un'incursione in un campo di cattivi, quando il
nostro tenente aveva ricevuto una chiamata a dieci minuti dalla zona di
atterraggio. I nemici avevano raggiunto un accordo e si erano ritrovati
all'improvviso dalla nostra parte. La nostra missione allora era
diventata proteggere i nostri nuovi amici da una forza di altri cattivi
che sgattaiolavano nella giungla per ucciderli. A quanto pare, il primo
gruppo di cattivi era stato dichiarato traditore dagli altri fuori di
testa ignoranti, perché non era composto da pazzi convinti come il più
pazzo dei fuori di testa. Quindi eravamo atterrati e avevamo fissato un
perimetro per proteggere i tizi cui avevamo sparato il giorno prima. Va
da sé che l'intera faccenda era in realtà una disputa tribale, ed
entrambe le parti avevano cercato di tenderci un'imboscata. Dopo essere
evasi, eravamo stati costretti a chiamare l'Aeronautica per risolvere la
situazione in modo diplomatico, con napalm, bombe a grappolo e
termobariche. Quando le cose cambiano, anche in modo radicale, guardi i
tuoi amici, scuoti la testa, alzi le spalle e ti adatti. Questo è quello
che fai, da soldato. I civili si arrabbiano quando da McDonald gli
cambiano il menu.

«Non fraintendermi, ci sono un sacco di specie che non hanno la capacità
di elaborare nuovi fatti e adattarsi, ma tra gli esseri umani accade per
lo più agli anziani, perché i loro cervelli non sono più in grado di
elaborare nuove informazioni. O sono solo pigri. Ma in alcune specie
anche i giovani hanno difficoltà ad accettare fatti che contrastano con
i loro rigidi sistemi di credenze. Sono colpito dagli umani. Ecco, l'ho
detto.»

«Quindi non siamo solo batteri?»

«Non montarti la testa. Ora siete batteri con del potenziale.»

\hfill\break

Trentasette minuti dopo, nella cabina di pilotaggio suonò un allarme.
«Eccellente!», riferì Skippy. «Due navi kristang sono saltate dentro,
una fregata che ci sta inviando segnali con un radiofaro e un
cacciatorpediniere.»

«Dove?», chiesi con ansia, mentre Skippy mostrava nello schermo sulla
paratia la nostra posizione rispetto a Paradiso, le due navi kristang e
diverse navi ruhar.

«Ho caricato una rotta nel sistema di navigazione.»

«Capito, signore», confermò Desai. «Reggetevi forte, pilota automatico
inserito.»

Mi ressi forte, mentre i motori del Dodo rombavano e scattavamo in
avanti. «Skippy, quanto dista? Quanto tempo manca all'incontro?»,
ripetei, con un occhio ansioso rivolto allo schermo. C'erano una mezza
dozzina di navi ruhar vicine in modo inquietante.

«La fregata kristang sta accelerando per intercettare la nostra rotta,
ci incontreremo fra tre minuti e 6,36749 secondi. Più o meno.»

Più o meno? Adams e io ci scambiammo un'occhiata divertita.

«Uhm», Skippy fece un'ottima imitazione di un grugnito, «il pilota
lucertola ha fatto una buona navigazione, sono saltati dentro vicino,
quasi quanto la loro tecnologia permette. Dev'essere stata fortuna,
nessuna viscida lucertola è così intelligente. Colonnello Joe, manca
poco, due navi ruhar si stanno preparando a saltare a breve raggio per
intercettare la fregata kristang. Non ho potuto occultare le lucertole
in salto, o i ruhar si sarebbero accorti che c'è qualcosa che non va nei
loro sistemi di sensori.»

«Ci vedono?»

«No, stanno tracciando la fregata. Il cacciatorpediniere kristang si sta
muovendo per fornirci copertura.»

«Grande. E sei sicuro che quando saremo a bordo della fregata e salterà
via riuscirai a fare spoofing sul suo sistema di navigazione, così
salteremo verso un posto diverso da quello del cacciatorpediniere?»

«Cosa? No, te l'ho detto, ho bisogno di essere fisicamente connesso
prima di poter stabilire qualsiasi genere di controllo sulla nave. Ma tu
pensa. Non mi stavi ascoltando?»

Porca puttana! «Skippy, che cazzo significa? Hai detto che pote\ldots»

«Ho detto che posso indurre qualsiasi nave ci prenda a bordo a saltare
in un posto diverso da qualunque altra nave la stia scortando.»

«Sì, e? Non hai bisogno di\ldots»

«Ah ah ah! Oh, sei carino, colonnello Joe. Pensi che abbia bisogno di
hackerare il computer delle lucertole per incasinare il loro motore di
salto? Non ci credo, amico! Distorcerò lo spazio-tempo all'ultimo
picosecondo, così il campo del loro motore di salto sarà dirottato.»

«Puoi distorcere lo spazio-tempo?», chiese Simms incredula.

«Signorsì», rispose Skippy con leggerezza. «È una specie di hobby. Ho
provato a collezionare francobolli, ma incasinare l'universo è tanto più
rilassante.»

Simms sollevò un sopracciglio e mi gettò uno sguardo eloquente,
dicendomi in silenzio come adesso capisse perché il piccolo pezzo di
merda l'avevo chiamato Skippy.

\hfill\break

La fregata si avvicinò sullo schermo, finché di colpo non fu proprio lì,
incombente sopra di noi, con una porta dell'hangar d'attracco già
aperta. Skippy lasciò che la fregata prendesse il controllo del sistema
di navigazione del Dodo e ci guidasse a bordo, la porta dell'hangar
aveva solo iniziato a chiudersi quando la fregata saltò via. «Bisogna
che ci muoviamo in fretta», esortò Skippy, «abbiamo appena saltato e ci
ho mandato fuori rotta, ora sto impedendo che si formi un campo di
salto, i kristang sanno che qualcosa non va, per ora pensano che il
problema sia il loro motore di salto, non sarà così per molto però.»

Le porte dell'hangar si chiusero con la lentezza di un'agonia e ci fu un
rombo quando l'hangar si ripressurizzò. Non appena l'indicatore di
pressione dell'aria raggiunse l'80\%, aprimmo la porta laterale e la
rampa posteriore. Mi scoppiavano le orecchie e sentivo un dolore
lancinante, che era difficile da ignorare. «La fotocamera è sotto il mio
controllo. Sta funzionando», disse Skippy, «si stanno bevendo la nostra
storia. Adesso la porta esterna è aperta.»

Skippy prese il controllo dei sensori di luce nella telecamera
dell'hangar d'attracco e fornì una falsa immagine alla nave; stava anche
chattando alla radio con l'equipaggio al ponte di comando. Con ``la
nostra storia'' intendeva dire che la falsa immagine che l'equipaggio
kristang vedeva attraverso la telecamera era eccitante e allettante: una
squadra del team delle forze speciali kristang che usciva da un Dodo
ruhar catturato con un oggetto molto speciale, di antica tecnologia
degli Anziani, che si presumeva fosse stato trovato su Paradiso. In base
a quanto deciso da Skippy, si trattava di un rubinetto a bolle di
energia, un dispositivo che estraeva energia libera dalle fluttuazioni
nella schiuma quantistica, o qualcosa del genere. «È una tecnologia
grezza, pensatela come una batteria che non si consuma mai. Credetemi,
questo impressionerà le stupide lucertole», aveva spiegato Skippy. In
ogni caso, funzionò: l'equipaggio del ponte di comando aprì la porta
all'interno della nave e attraversammo l'hangar a gravità zero, solo
poche persone mancarono il bersaglio e furono tirate dentro prima che
rimbalzassero contro il muro e galleggiassero via. Poiché le fregate
kristang avevano equipaggi ridotti e Skippy determinò che questa fregata
era a corto di due membri dell'equipaggio regolamentare, non c'era
nessuno di loro all'hangar d'attracco quando arrivammo. L'equipaggio del
ponte di comando disse a Skippy che stavano mandando qualcuno a
prenderci, il resto dell'equipaggio era occupato a cercare febbrilmente
di capire perché il loro motore di salto non funzionasse.

Giraud, in testa, riuscì ad aprire la porta esterna della camera
d'equilibrio, non bisognava fare altro che premere un pulsante e poi
tirare una leva. All'interno della camera d'equilibrio, tenni Skippy in
una mano e inserii la spina in una presa, mentre Giraud digitava nel
pannello di controllo una sequenza che la nostra lattina di birra
super-intelligente gli aveva insegnato, per forzare l'apertura della
porta interna mentre quella esterna non era ancora chiusa. Era un
momento pericoloso: con entrambe le porte aperte, avrebbe suonato un
allarme che neppure Skippy avrebbe potuto silenziare. La porta interna
si aprì con un botto e un forte allarme prese a rimbombare; le ventidue
persone della squadra d'assalto si ammassarono per attraversare la
camera d'equilibrio e uscire in corridoio il più in fretta possibile per
degli umani e, non appena l'ultimo di loro fu passato, colpii con un
pugno il pulsante che fece chiudere con fragore la porta esterna e
interruppe l'allarme. Era il segnale per Desai di far muovere il Dodo e,
per la nostra squadra d'assalto, di dividersi e dirigersi verso il ponte
di comando e la sala macchine.

Il compito più importante era il mio: spinottare Skippy in una presa a
muro. Una presa a muro che, per un'impanicata eterna frazione di
secondo, non riuscivo a trovare. In mia difesa, c'è da dire che stavo
galleggiando a testa in giù vicino al soffitto, cercando di tenermi
lontano dai componenti della squadra d'assalto, prendendomi urti,
gomitate e anche un calcio in faccia. Eravamo tutti impacciati a gravità
zero, nessuno di noi si era allenato per il combattimento in queste
condizioni e non avevamo avuto opportunità di fare pratica. A bordo del
Dodo, avevo costretto tutti, tranne i piloti, a esercitarsi nel
volteggio, nell'atterraggio controllato su un muro, spingendo con piedi
e mani. Non c'era abbastanza spazio, o tempo sufficiente nel Dodo,
perché tutti acquisissero una confidenza tale da poter fare molto di più
che evitare di vomitare mentre vorticavano in giro.

Alla fine notai la presa a muro, nello stesso momento risuonò il rumore
di armi da fuoco davanti a noi: i kristang, diretti all'hangar
d'attracco, avevano visto l'avanguardia della squadra d'assalto. Dal
suono, riconobbi il ronzio delle armi ruhar e il più pesante rimbombo di
un fucile kristang, poi solo ronzio. Ignorando le distrazioni, mi spinsi
giù dal muro, presi la spina che tenevo tra i denti e la inserii con
cura nella presa. «Ci sei?»

«Occupato», fu l'unica risposta di Skippy e, considerando la sua
fulminea velocità di elaborazione, mi preoccupai. Poi: «Ci sono. Siamo a
posto. Ho sotto controllo i sistemi della nave. Spingo le porte
dell'hangar d'attracco». Ci fu una vibrazione, poiché Skippy aveva usato
una procedura d'emergenza per far saltare le porte verso l'esterno,
piuttosto che ritrarle verso i lati, non potevamo aspettare che
eseguissero il loro ciclo normale. «Il Dodo si sta muovendo.»

La nave sbandò con violenza, facendomi rimbalzare via dal muro. Skippy
aveva detto che aveva il controllo, la nave non avrebbe dovuto essere in
grado di muoversi! «Che cazzo era?»

«Il Dodo ha colpito la porta dell'hangar d'attracco. Il Dodo ha
riportato un danno sostanziale. È funzionale alla missione.»

Desai aveva preso male le misure dell'uscita in un veicolo sconosciuto e
in una situazione complicata dall'aria che fuoriusciva dall'hangar
d'attracco. Skippy riteneva che il Dodo potesse ancora portare a termine
la sua missione, che era quella di volare davanti alla fregata e far
saltare il ponte di comando della nave con i cannoni del Dodo. I
kristang, che erano da tempo abituati alla pirateria tra clan e
all'interno di clan, avevano ideato le loro navi contro chi andava
all'arrembaggio come noi, anche se i progettisti avevano pensato alle
fazioni rivali dei kristang, non ai primitivi umani. La porta del ponte
di comando era blindata, abbastanza da far sì che sfondarla avrebbe
provocato una falla nella nave ed esposto la nostra gente al vuoto.
Finché i kristang erano dietro quella porta, potevano impedirci di usare
il mezzo, potevano anche disabilitarlo o autodistruggerlo, a prescindere
dal fatto che Skippy fosse inserito. Era quasi impossibile impadronirsi
del ponte di comando di una nave da guerra kristang.

Perciò non ce ne saremmo impadroniti, l'avremmo fatta a pezzi
dall'esterno. Se il Dodo fosse stato ancora funzionante, se Desai fosse
stata in grado di pilotarlo con efficacia senza quasi nessuna esperienza
e se fosse stata in grado di controllare i cannoni del Dodo al punto da
colpire il ponte di comando senza far saltare in aria il resto della
nave. Se. Era tutto un grandissimo se.

Ci furono spari sul davanti e a poppa, la gente urlava e urlava;
qualcuno lanciò una granata stordente, poi la nave vibrò di nuovo con
violenza. E ancora e ancora. «Il capitano Desai è riuscita a distruggere
il ponte di comando. Anche parte della prua della nave e dieci metri del
lato di dritta a poppa del ponte.»

«Merda! Il danno è così grave?»

«Ha fatto bene per la sua mancanza di esperienza. L'ulteriore danno sarà
la prova convincente che questa nave è stata colpita da una mina stealth
ruhar e non influenzerà le funzioni del velivolo utili al nostro scopo.»

È stato detto che un comandante ha il lavoro più duro in combattimento
perché, dopo che le persone sono state addestrate e i piani sono stati
elaborati, deve sedersi e stare a guardare, incapace di fare molto per
influenzare il risultato, mentre la sua gente lotta. Nella mia
esperienza molto limitata, quel detto è una stronzata al 100\%. Il
lavoro più duro non è quello del comandante, il lavoro più duro è svolto
dai soldati semplici che imbracciano fucili, esponendosi al fuoco
nemico. Sono i soldati semplici a combattere e morire. Il lavoro più
duro è il loro. Quello del comandante è il lavoro più solitario, il
lavoro che ti fa sentire in colpa e inutile per il fatto di non essere
in prima linea a fare qualcosa di utile per i tuoi amici. A
combattimento finito, ho sempre pensato che le cose sarebbero andate
meglio se fossi stato lì, che forse avrei visto il nemico e gridato un
avvertimento prima che qualcuno venisse colpito; forse avrei ucciso il
nemico prima che prendesse uno dei nostri. Forse avrei potuto migliorare
le cose. O forse sarei stato soltanto ucciso. In ogni caso, avrei fatto
qualcosa. L'unica cosa che feci, dopo aver inserito Skippy, fu tenermi
stretto a un muro e cercare di seguire le battaglie a prua e a poppa sul
mio zPhone. Era più che caotico, non riuscii a farmi un'idea chiara di
quello che stava accadendo fino a quando la lotta non fu finita. Alla
faccia del mio comandare qualcosa.

Le nostre perdite nella presa della nave furono quattro, e ci furono tre
feriti gravi. Quattro morti e tre feriti su ventidue delle squadre di
assalto. Un terzo della nostra forza era ormai andato, oppure
impossibilitato a combattere. In una battaglia. Nel prendere una nave.

La nave l'avevamo presa, era nostra. Giraud ebbe solo una vittima nel
raggiungere il centro di controllo della nave: il soldato Arun Kurien
dell'esercito indiano morì per essere stato colpito da uno dei due
kristang che stavano bloccando loro l'accesso. Tutte le altre vittime
erano nella squadra di Chang, che aveva sequestrato la sala macchine.
Avevano fatto il lavoro più duro, dovevano uccidere cinque kristang che
erano disperatamente motivati a impedire che la loro nave venisse presa.
Tre delle lucertole erano state uccise nello scontro a fuoco iniziale,
quando eravamo in vantaggio in virtù dell'effetto sorpresa, con le
restanti due fu molto più difficile. Una di loro era riuscita a
procurarsi una tuta in armatura potenziata, nonostante Skippy avesse
bloccato la porta dell'armadietto dove le custodivano; il kristang aveva
in parte indossato e quasi acceso l'armatura, quando tre dei nostri
avevano sparato alla cieca nel compartimento per fornire a Chang la
copertura per buttare una granata. Il kristang era tosto, stava ancora
cercando di abbottonare la tuta dopo l'impatto della granata e ci era
voluto un fuoco concentrato per farlo fuori.

Il problema peggiore era stato l'ultimo kristang. Doveva avere capito di
essere l'unico rimasto, nonostante Skippy avesse spento il sistema di
comunicazione, e aveva deciso che avrebbe fatto saltare in aria la nave,
piuttosto che arrendersi. Skippy aveva avvertito Chang che doveva
fermare quel kristang subito, subito, subito, subito, prima che fosse in
grado di penetrare nel contenimento del reattore. Per quanto Skippy
facesse tutto il possibile, i controlli della scarica di plasma del
reattore erano manuali, nulla su cui potesse interferire nel tempo
disponibile.

Chang mi aveva detto quello che era successo; la squadra d'assalto aveva
già perso due persone, uccise nel tentativo di raggiungere il kristang
nello spazio ristretto in cui si stava nascondendo. Vedendo che la
situazione era disperata, il sergente Yu Qishan aveva tolto la sicura a
due granate e si era lanciato dietro l'angolo, morendo quando il
kristang gli aveva sparato alla testa. Anche quell'ultima lucertola era
morta quando le granate erano esplose e avevano aperto una falla nello
scafo della nave, risucchiando sia lui che il sergente Yu nello spazio.
Rischiando di risucchiare la maggior parte della squadra d'assalto di
Chang nello spazio, non fosse che una paratia si era abbassata con
violenza in automatico per evitare ulteriori perdite d'aria nella sala
macchine.

Chang cercò davvero di consolarmi, anche se il sergente Yu era uno dei
suoi uomini. «Sapeva che la sopravvivenza dell'umanità dipendeva
dall'impedire che quel kristang distruggesse il reattore. Ha fatto il
suo dovere. Gli renderemo onore quando torneremo a casa.»

Aveva ragione. Continuavo a sentirmi una merda.

\hfill\break

All'indomani della battaglia, Skippy aprì la porta dell'hangar
d'attracco di babordo per poter imbarcare il Dodo danneggiato, e una
Desai ancora turbata raggiunse subito il centro di controllo della nave
per fare un corso accelerato di pilotaggio di una fregata kristang. Il
mio primo istinto fu quello di dare uno stop, per lasciare che tutti si
riprendessero e, soprattutto, ci si occupasse dei feriti. Simms e Skippy
mi esortarono a continuare con il piano: subito, non c'era un minuto da
perdere. Simms mi assicurò che si stava prendendo cura dei feriti, che
erano stati portati nella piccola infermeria della fregata, e non c'era
niente che potessi fare per aiutarli. Con le misere forniture mediche
che avevamo portato con noi, nessuno avrebbe potuto fare molto per loro.
La cosa migliore da parte mia, disse Skippy, era continuare con il piano
e prendere una nave madre thuranin, dove i feriti avrebbero potuto
essere curati con la tecnologia medica incredibilmente avanzata dei
cyborg. Le ferite avrebbero potuto guarire, perfino gli arti avrebbero
potuto ricrescere, mi assicurò Skippy, una volta che messi i nostri
soldati feriti nelle vasche curative dei thuranin. Mi sembrava
sbagliato. Ingoiai il mio orgoglio e lasciai che la mia gente facesse il
suo lavoro.

Mentre Desai premeva pulsanti sotto la direzione di Skippy, programmando
un salto, io mi guardavo attorno nel centro di controllo con stupore.
Avevamo una nave. Un'astronave. Noi. Ignobili, ignoranti scimmie
terrestri a bassa tecnologia. Un'astronave.

Cazzo.